% !PRINTED 266
\weberpage{266}
%<<CONT
sein müsste, und es würden $a$ und $d$ nicht reell ausfallen. Also
nehmen wir für $a$ und $c$, wie oben, die $p\,(p-1)$ verschiedenen
Werthsysteme, und erhalten aus (11) für jedes von ihnen
$1/2\,(p-1)$ Bestimmungen von $b$ und $d$; auch diese Zahl ist zu
halbiren und es ergeben sich
\[
1/4\;p\,(p-1)\,(p-1)
\]
Substitutionen der III\textsuperscript{ten} Art. Die Gesammtzahl aller Substitutionen
 der Gruppe $L_p$ ist hiernach, wie es sein muss,
\[
p^2 + 1/4\,p\,(p-3)\,(p+1) + 1/4\,p\,(p-1)^2 = 1/2\,p\,(p^2-1).
\]

\begin{center}
§. 69.

\gesperrt{Imaginäre Form der Gruppe $L_p$.}
\end{center}

Die Substitutionen der beiden ersten Arten der Gruppe $L_p$
haben die Eigenschaft, dass die ihnen entsprechende Normalform
 $T$ oder $S$ gleichfalls in $L_p$ enthalten ist, und dass man sie
in diese Normalformen transformiren kann durch \gesperrt{reelle} Substitutionen
, d. h. durch Substitutionen, die selbst in $L_p$ vorkommen
. Beides trifft bei den Substitutionen der dritten Art
nicht zu.

Man kann aber, wie wir jetzt beweisen werden, die ganze
Gruppe $L_p$ in eine andere Gruppe $\Lambda_p$ so transformiren, dass $\Lambda_p$ ein
mit $L_p$ \gesperrt{isomorpher Theiler der Gruppe} $E_{2,p}$ ist, dass die
Normalformen der Transformationen dritter Art in $\Lambda_p$ enthalten
sind und dass jede Substitution dritter Art in $\Lambda_p$ in die
Normalform transformirt werden kann durch Substitutionen von
$\Lambda_p$ selbst.

Um eine solche Transformation von $L_p$ zu finden, betrachten
wir die Substitution
\weq{1}{S = \begin{pmatrix} \sigma, & 0 \\ 0, & \sigma^{-1} \end{pmatrix}}
unter der Voraussetzung, dass $\sigma$ und $\sigma^{-1}$ conjugirt imaginär
sind, und suchen die Substitution
\weq{2}{R = \begin{pmatrix} \lambda, & \mu \\ \nu, & \varrho \end{pmatrix}}
so zu bestimmen, dass
\weq{3}{R\,S\,R^{-1}}
%CONT>>
